\statement{A.76}{Proved}{An actual finite central refinement does not suffice for global-negative framed order preservation}
There are a finite actual quotient \(Q=\Pi/N\), its actual
source-central refinement
\[
 Q^\dagger=\Pi/\overline{[N,\Pi]},
 \qquad \Pi=C_2\amalg C_3,
 \tag{A.76.1}
\]
and positive peripheral data on \(Q^\dagger\), fixing the literal
cusp, with the following properties. Their negative companion has
order two and has prescribed s- and r-normalizers of order two.
An ordinary full positive framed lift exists. Nevertheless no full
positive peripheral h fixing \(\tau\) literally, inducing the
specified finite action, and retaining the specified finite s-frame
can satisfy
\[
                         (hc)^2=1.
 \tag{A.76.2}
\]
Thus the failure of framed order-preserving lifting is present for
the globally negative multiplier \(-1\), after the actual finite
central-refinement test and both quotient normalizer norms have
passed. It does not exclude reselection of the finite frame or
refute the unframed existence problem.

\begin{proof}
\emph{The original actual data and finite central refinement.}\par
Use the original A.56 witness, before the subsequent affine changes
in A.64 and A.66. Write \(D=\Pi/U\) for its actual quotient over C.
By A.56,
\[
 U=\overline{[U,\Pi]},
\]
and D has a positive peripheral action h, a negative involution
\(b=hc\), and literal s- and r-frames with both negative norms one.
The ordinary full positive framed lift exists by A.3 and A.56.

Its finite old model \(D_0\) has normal forms \((a,m,q)\),
with \(a\in V^8\), \(m\in\operatorname{Sym}^2_k(k^8)\),
and \(q\in H\), using exactly the opposite-weight law and
C4 carry in A.56. The finite head and source are (A.56.2)--(A.56.4);
in the original cocycle coordinates the source elements are
\[
 s\longmapsto(d(s),0,H_s),\qquad
 r\longmapsto(0,0,H_r).
\]
The actual polynomial positive and canonical actions are those of
\path{work_A/check_Q6_central_cover.py}, with the original saved
corrections. In particular \(D_0\) is an actual, action-stable
projection of D, of order
\[
                         |D_0|=2^{972}\,72\,7^{13}.
 \tag{A.76.3}
\]
Take \(Q=D_0\) and \(N=\ker(\Pi\to D_0)\). Since \(U\subseteq N\),
\[
 U=\overline{[U,\Pi]}\subseteq\overline{[N,\Pi]}.
 \tag{A.76.4}
\]
Therefore the actual refinement (A.76.1) is a quotient of D.
It inherits the negative involution and both norm-one frames,
not merely corresponding data on a conveniently chosen cover.
The original full framed lift also realizes these descended data.

The refinement is finite. Indeed its central subgroup
\(Z=N/\overline{[N,\Pi]}\) has finite index \(m=|Q|\)
and is topologically finitely generated by the finite-index
Schreier argument. The continuous transfer to Z restricts on Z
to \(z\mapsto z^m\). Its image has exponent dividing six,
because the source generators have orders two and three.
Consequently Z has exponent dividing \(6m\). A finitely generated
abelian profinite group of bounded exponent is finite, proving
the assertion. This also gives a direct finite construction;
no unspecified limiting quotient is used in (A.76.1).

The old witness is
\path{work_A/Q6_actual_central_solution.json}.
Its SHA-256 begins \texttt{b042bf92e479019c}; the complete hash,
together with those of all obstruction inputs, is recorded in
\path{work_A/A76_original_witness_hashes.json}.
In particular this theorem concerns the original stored datum,
not one of the later repaired tuples.

\emph{Every candidate preserves the actual six-cup quotient.}\par
Put \(T_0=V^{(2)}\), interpreted as the underlying binary
H-module with its displayed squared monomial action. The finite
old residual \(P_0\) satisfies \(P_0^{\rm ab}=V^8\).
The normal odd subgroup of H has no invariants on \(T_0\), so
coprime Hochschild--Serre gives
\[
 H^1(H,T_0)=H^2(H,T_0)=0,
 \qquad H^1(D_0,T_0)
       \simeq\operatorname{Hom}_H(P_0^{\rm ab},T_0)
       \simeq k^8.
 \tag{A.76.5}
\]
Here the Hom is over the underlying additive binary modules;
the Frobenius comparison supplies the displayed k-dimension.
The source cocycles are the complete eight classes, so inflation
in (A.76.5) is an isomorphism onto \(H^1(\Pi,T_0)\).
The five-term sequence therefore makes the actual transgression
\[
 \operatorname{Hom}_{D_0}(N,T_0)
              \longrightarrow H^2(D_0,T_0)
 \tag{A.76.6}
\]
injective.

The six literal cusp cups \(e_i\smile e_j\), \(0\leq i<j\leq3\),
are the independent actual extension classes of A.63. Their finite
source quotient is
\[
 D_0\mathbin{\times}_{V^4\rtimes H}(N_6\rtimes H),
 \tag{A.76.7}
\]
where \(N_6\) has the coordinatewise cross-cup law and kernel
\(T_0^6\). Its actual source attainment is the same finite
Frattini argument as A.63. The retained old positive action fixes
the four cusp cohomology classes, hence all six cup classes.
Equivariance and the injectivity in (A.76.6) force preservation
of their six actual kernel maps and force identity on their
coefficient quotient. Canonical c similarly has its prescribed
squared monomial action. Thus every proposed full lift of the
retained data descends to (A.76.7); kernel stability and coefficient
action are conclusions, rather than extra restrictions on candidates.

\emph{Exhausting all source cochains and frame changes.}\par
Two descended positive actions with the same old \(D_0\)-action
and coefficient action differ by a \(T_0^6\)-valued cocycle on
\(D_0\). By (A.76.5), its source inflation ranges over exactly
the complete 44-dimensional cocycle space in each coefficient
copy: eight cohomology classes and 36 coboundaries. The literal
cusp condition kills its cusp restriction. Keeping the same old
s-frame makes its s-restriction a coboundary. The actual joint
restriction isomorphism (A.56.5), with the Frobenius comparison,
therefore forces the difference to be
\[
 \operatorname{Inn}(a),\qquad a\in(T_0^6)^\tau.
 \tag{A.76.8}
\]
The negative coefficient action is identity on this cusp-fixed
space. If \(b'=\operatorname{Inn}(a)b\) and the frame is
transported to \(\beta'=a\beta\), direct multiplication gives
\[
 \beta'^{-1}b'(\beta'^{-1})
                 =\beta^{-1}b(\beta^{-1}).
 \tag{A.76.9}
\]
Every remaining frame over the same old frame differs by an element
of \((T_0^6)^s\), and changes the square by its complete
\(1+D_s\) norm. Old globally central frame changes have no new
\(T_0\)-component in (A.76.7). Thus (A.76.8)--(A.76.9) exhaust
all admissible source cochains and all new frame variations for
this prescribed old datum.

\emph{The exact nonzero invariant.}\par
In each coefficient copy, \((T_0)^s\) has k-dimension 20;
the frame norm has rank eight and invariant target dimension 12.
At zero-based physical weight 23, the frame coefficient involution
fixes the weight literally with scalar one. Evaluation there
annihilates the entire norm image. In cup order
\((01,02,03,12,13,23)\), the original-law s-square has values
\[
                           (0,0,5,0,5,5)\in k^6.
 \tag{A.76.10}
\]
The value five in the 03 coordinate is nonzero. It therefore
cannot be removed by any of the variations just exhausted.

The independent direct certificate
\path{work_CE/code/verify_A63_s_norm_direct.py}
checks all 360 binary s-centralizer norm columns, all 216
coefficient-action columns, and all 72 literal-cusp inner
correction columns. The complete results are
\path{work_CE/A63_s_norm_independent_certificate.json}
and \path{work_CE/review_A63_fixed_affine_s_norm.txt}.
Their hashes are bound in the witness record above. This is an
exact norm-image annihilator, not a search timeout or a statement
about one unvaried representative.

\emph{Why a full negative involution would contradict it.}\par
Suppose a full positive peripheral h fixes \(\tau\) literally,
retains the specified finite action and s-frame, and has
\(b=hc\) involutory. Let \(\beta\) be its full s-conjugator
with that prescribed projection. Since
\(b(s)=\beta s\beta^{-1}\), involutivity gives
\[
              \beta^{-1}b(\beta^{-1})\in C_\Pi(s)=\langle s\rangle.
 \tag{A.76.11}
\]
The free-factor centralizer equality is the full profinite
free-product theorem.\cite{A_HR1985}
The retained S3 quotient has trivial b-action, and \(\beta\)
has trivial image there. Hence the element in (A.76.11) maps to
one in S3, while s has nontrivial image. It must therefore be
one in \(\Pi\). Its image in the actual quotient (A.76.7) would
be a norm-one s-frame, contradicting (A.76.10).

The multiplier of b in this argument is the global unit \(-1\);
the positive lift has multiplier one. The literal cusp and the
specified finite frame are part of the lifting problem. No claim
is made for unnormalized lifts with only an unrelated quotient
cusp framing, or for existential reselection of the old finite
affine datum. The later constructions A.64 and A.75 change that
datum and remain fully compatible with this counterexample.
\end{proof}
